\documentclass[a4paper,12pt]{article}
\usepackage[utf8]{inputenc}
\usepackage[T1]{fontenc}
\usepackage[french]{babel}
\usepackage{geometry}
\geometry{margin=2.5cm}

\title{Cybersécurité - Cours}
\author{Mignard Mael}
\date{\today}

\begin{document}

\maketitle

\section{Lexique de la cybersécurité}
\begin{itemize}
    \item \textbf{Cybersécurité} : Ensemble des pratiques, technologies et processus visant à protéger les systèmes informatiques, les réseaux et les données contre les attaques, les dommages ou les accès non autorisés.\\
    \item \textbf{Cyberespace} : Environnement virtuel créé par les réseaux informatiques où les interactions humaines et les échanges d'informations ont lieu.\\
    \item \textbf{SI} : Systeme d'information, ensemble des ressources destinées à collecter, stocker, traiter et diffuser l'information au sein d'une organisation.\\
    \begin{itemize}
        \item \textbf{Actif} : Tout élément d'information, matériel ou logiciel, qui a de la valeur pour l'organisation et qui doit être protégé.
        \item \textbf{Vulnérabilité} : Faiblesse ou faille dans un système qui peut être exploitée par une menace pour compromettre la sécurité.
        \item \textbf{Menace} : Potentiel danger qui peut exploiter une vulnérabilité pour causer un dommage.
        \item DI
    \end{itemize}
\end{itemize}


\end{document}
